\documentclass[a4paper,12pt]{article}

\usepackage[T2A]{fontenc}
\usepackage[utf8]{inputenc}
\usepackage[russian]{babel}
\usepackage{amsmath,amssymb,amsthm,mathtools}
\usepackage{helvet}
\renewcommand{\familydefault}{\sfdefault}
\usepackage{icomma}
\usepackage{geometry}
\usepackage{microtype}
\usepackage{tikz}
\usetikzlibrary{calc,arrows.meta,positioning,shapes.geometric}
\usepackage{tabularx,booktabs,longtable}
\usepackage{enumitem}
\usepackage{hyperref}
\usepackage{parskip}
\usepackage{fancyhdr}
\usepackage{setspace}
\usepackage{xcolor}

\geometry{top=20mm,bottom=20mm,left=22mm,right=22mm}
\onehalfspacing

\definecolor{accent}{HTML}{008080}
\definecolor{darkaccent}{HTML}{004D4D}
\definecolor{textgray}{HTML}{333333}

\pagestyle{fancy}
\fancyhf{}
\fancyhead[L]{Производная сложной функции}
\fancyhead[R]{Николь Саркисьянц}
\fancyfoot[C]{\thepage}

\tikzset{
startstop/.style={
ellipse,
draw=accent,
minimum width=3.5cm,
minimum height=0.8cm,
align=center
},
process/.style={
rectangle,
rounded corners,
draw=accent,
minimum width=5cm,
minimum height=0.9cm,
align=center
},
decision/.style={
diamond,
draw=accent,
aspect=2,
minimum width=3.5cm,
align=center
},
arrow/.style={
-{Latex[length=3mm]},
thick
}
}

\begin{document}

\begin{titlepage}
\centering

\vspace*{2cm}

{\Large ЕГЭ по профильной математике\par}

\vspace{1cm}

{\huge\bfseries Чек-лист\par}

\vspace{0.5cm}

{\Large Производная сложной функции\\
Константы и переменные в производных}

\vspace{2cm}

{\large
Николь Саркисьянц
}

\vfill

{\large
Лёвин Артём Александрович\\
эксклюзивно для Николь Саркисьянц
}

\vspace{1cm}

{\large \today}

\vspace{1.5cm}

\begin{tikzpicture}
\draw[color=accent,line width=1.2pt]
(0,0) .. controls (2,1.2) and (4,-1.2) .. (6,0)
.. controls (8,1.2) and (10,-1.2) .. (12,0);
\end{tikzpicture}

\end{titlepage}

\section*{1. Главная идея производной сложной функции}

Во многих задачах встречается конструкция

\[
y=f(g(x)).
\]

Внутри находится \textbf{внутренняя функция} $g(x)$, снаружи --- \textbf{внешняя функция} $f$.

Правило:

\[
\boxed{(f(g(x)))'=f'(g(x))\cdot g'(x)}
\]

То есть:

\begin{enumerate}
\item найти внутреннюю функцию;
\item вычислить её производную;
\item вычислить производную внешней функции;
\item перемножить результаты.
\end{enumerate}

\subsection*{Пример 1}

\[
y=\sqrt{x^2-x-6}
\]

Вводим замену:

\[
m=x^2-x-6.
\]

Тогда

\[
(m(x))'=2x-1.
\]

Внешняя функция:

\[
y=\sqrt m.
\]

По формуле:

\[
(\sqrt m)'=\frac1{2\sqrt m}.
\]

Возвращаем замену:

\[
\boxed{
y'=\frac{2x-1}{2\sqrt{x^2-x-6}}
}
\]

\textbf{ОДЗ:}

\[
x^2-x-6>0.
\]

\[
(x-3)(x+2)>0,
\]

следовательно,

\[
x<-2\quad\text{или}\quad x>3.
\]

\section*{2. Производная экспоненты со сложной степенью}

Если

\[
y=e^{g(x)},
\]

то

\[
\boxed{
y'=e^{g(x)}\cdot g'(x)
}
\]

Важно:

производится вся степень.

Пример:

\[
y=e^{5x^2}
\]

Степень:

\[
g(x)=5x^2.
\]

Производная степени:

\[
g'(x)=10x.
\]

Ответ:

\[
\boxed{
y'=10xe^{5x^2}
}
\]

\section*{3. Константы и переменные}

Перед началом решения необходимо определить:

\[
\boxed{\text{Что является переменной?}}
\]

Числа и параметры считаются константами.

Например:

\[
y=a x^2+b x+c,
\]

где $a,b,c$ --- константы.

Тогда:

\[
y'=2ax+b.
\]

Параметры не дифференцируются.

\subsection*{Особый случай}

Если написано:

\[
y(a)=a^2\sin x+3b\sin x,
\]

то переменной является $a$.

Следовательно:

\[
\sin x
\]

становится константой.

Тогда:

\[
(y(a))'
=
2a\sin x.
\]

\section*{4. Производная степенной сложной функции}

Для

\[
y=(g(x))^n
\]

используется:

\[
\boxed{
y'=n(g(x))^{n-1}g'(x)
}
\]

Пример:

\[
y=(x^2-x-6)^{10}.
\]

Внутренняя функция:

\[
g(x)=x^2-x-6.
\]

Тогда:

\[
g'(x)=2x-1.
\]

Получаем:

\[
\boxed{
y'=10(x^2-x-6)^9(2x-1)
}
\]

\clearpage



\section*{6. Мини-памятка}

\[
\boxed{
\text{Сначала найдите переменную.}
}
\]

\[
\boxed{
\text{Затем найдите внутреннюю функцию.}
}
\]

\[
\boxed{
\text{Производная сложной функции = внешняя производная}\times
\text{внутренняя производная.}
}
\]

\clearpage

\section*{Тренировочный блок}

Решите производные.

\begin{enumerate}
\item
\[
y=\sqrt{3x^2+5x-1}
\]

\item
\[
y=e^{4x^2-x}
\]

\item
\[
y=(2x^3-x+4)^5
\]

\item
\[
y=\sin(x^2+3x)
\]

\item
\[
y=(x^2-2x+1)^{12}
\]
\end{enumerate}

\clearpage

\section*{Ответы}

\renewcommand{\arraystretch}{1.35}

\begin{center}
\begin{tabularx}{\linewidth}{cX}
\toprule
№ & Ответ\\
\midrule
1 &
$\displaystyle
y'=\frac{6x+5}{2\sqrt{3x^2+5x-1}}
$\\

2 &
$\displaystyle
y'=(8x-1)e^{4x^2-x}
$\\

3 &
$\displaystyle
y'=5(2x^3-x+4)^4(6x^2-1)
$\\

4 &
$\displaystyle
y'=(2x+3)\cos(x^2+3x)
$\\

5 &
$\displaystyle
y'=12(x^2-2x+1)^{11}(2x-2)
$\\
\bottomrule
\end{tabularx}
\end{center}

\end{document}